\section{Support Vector Machines}

\acsp{svm} extend Support Vector Classifiers to the case of non-linear decision surfaces. There are two main ideas behind this extension:
\begin{enumerate}
    \item Augment the feature space so that observations become linearly separable in that new space.
    \item Can we actually do this without going in a higher dimensional space?
\end{enumerate}

\paragraph{Quadratic decision surface} Suppose we have $p$ features $x_1, \dots, x_p$. We can build the augmented feature space in this way:
\begin{equation}
    x_1, x_1^2, x_2, x_2^2, \dots, x_p, x_p^2.
\end{equation}
The \acs{svc} problem on the augmented space becomes:
\begin{equation}
    \begin{aligned}
         & \min_{\beta_0, \betavec, \epsilonvec} \frac{1}{2}\norm{\betavec}^2                                                                                        \\
         & \text{s.t.}\; y_i\left(\beta_0 + \sum_{j = 1}^{p}\beta_{j1}x_{ij} + \sum_{j = 1}^{p}\beta_{j2}x_{ij}^2\right) \geq 1 - \epsilon_i, \quad i = 1, \dots, n, \\
         & \phantom{s.t.} \epsilon_i \geq 0, \quad i = 1, \dots, n,                                                                                                  \\
         & \phantom{s.t.} \sum_{i=1}^n \epsilon_i \leq C.
    \end{aligned}
\end{equation}

In the augmented feature space $\rr^{2p}$ we obtain indeed a linear decision surface:
\begin{equation}
    \zvec \in \rr^{2p}\negthinspace\colon \beta_0 + \betavec^\t\zvec = 0.
\end{equation}
In the original feature space $\rr^p$, instead, we obtain a quadratic decision surface:
\begin{equation}
    \xvec \in \rr^p\negthinspace\colon \beta_0 + \sum_{j = 1}^{p}\beta_{j1}x_{j} + \sum_{j = 1}^{p}\beta_{j2}x_{j}^2 = 0.
\end{equation}
The space can be extended considering interaction terms $x_i x_j$ or even higher degrees, but it becomes computationally expensive.

\paragraph{Going higher dimensions}
Recall that in an \acs{svc} the decision surface is linear and defined by:
\begin{equation}
    \begin{aligned}
        f(\xvec) & = \beta_0 + \betavec^\t\xvec                                \\
                 & = \beta_0 + \sum_{j=1}^p\beta_jx_j                          \\
                 & = \beta_0 + \betavec \cdot \xvec                            \\
                 & = \beta_0 + \sum_{i = 1}^n \gamma_i y_i \xvec_i\cdot \xvec.
    \end{aligned}
\end{equation}
In the optimization (estimation of $\gammavec$), we solve:
\begin{equation}
    \max_{\gammavec} \sum_{i = 1}^n \gamma_i - \frac{1}{2}\sum_{i,j = 1}^{n} \gamma_i\gamma_j y_i y_j \xvec_i \cdot \xvec_j.
\end{equation}
We notice that the optimization depends on the data only through the inner products $\xvec_i \cdot \xvec_j$. Therefore, it would make sense to replace the inner product with a more general function $K(\xvec_i, \xvec_j)$, called \emph{kernel}.

Suppose we have a function $\Phi\colon \rr^p \to \rr^N$ with $N > p$, that maps the original feature space to a higher-dimensional space. Therefore, we can rewrite $f(\xvec)$ in this space as:
\begin{equation}
    f(\xvec) = \beta_0 + {\underbracket{\betavec}_{\mathclap{\betavec\,\in\,\rr^N}}}^\t \underbracket{\Phi(\xvec)}_{\zvec} = \beta_0 + \sum_{i = 1}^n\negthinspace\underbracket{\gamma_i}_{\gammavec\,\in\,\rr^n} \negthickspace\negthickspace y_i \Phi(\xvec_i) \cdot \Phi(\xvec).
\end{equation}
So, $\beta_0$ and $\gammavec$ get estimated by:
\begin{equation}
    \max_{\gammavec} \sum_{i = 1}^n \gamma_i - \frac{1}{2}\sum_{i,j = 1}^{n} \gamma_i\gamma_j y_i y_j \Phi(\xvec_i) \cdot \Phi(\xvec_j).
\end{equation}
\begin{note}
    We notice that while the $\beta$s are $N$, the $\gamma$s stay in number $n$.
\end{note}
Since $\Phi$ is generally not known, we would not want to compute $\Phi(\xvec_i) \cdot \Phi(\xvec_j)$ directly. However, we can rely on the so called \emph{kernel trick}, which consists in finding a function $K\negthinspace\colon \rr^p \times \rr^p \to \rr$ such that it acts as a dot product in some higher dimensional space for some function $\Phi$. Mercer's theorem states the properties that such function needs to have.

Using such kernel in $f(\xvec)$:
\begin{equation}
    f(\xvec) = \beta_0 + \sum_{i = 1}^n \gamma_i y_i K(\xvec_i, \xvec)
\end{equation}
and the optimization becomes:
\begin{equation}
    \max_{\gammavec} \sum_{i = 1}^n \gamma_i - \frac{1}{2}\sum_{i,j = 1}^{n} \gamma_i\gamma_j y_i y_j K(\xvec_i, \xvec_j).
\end{equation}
Therefore, the decision surface is defined by:
\begin{equation}
    \xvec\colon f(\xvec) = \beta_0 + \sum_{i=1}^n \gamma_i y_i K(\xvec_i, \xvec) = 0
\end{equation}
and it is non-linear, with the shape influenced by the choice of the kernel.

The most common choices for the kernel are:
\begin{description}[beginpenalty=10000]
    \item[Linear Kernel] $\hfill K(\xvec, \zvec) = \xvec \cdot \zvec. \inlineeqno$
    \item[Polynomial Kernel] $\hfill K(\xvec, \zvec) = (1 + \xvec^\t\zvec)^m \text{ of degree } m. \hfill \inlineeqno$
    \item[Radial Kernel] $\hfill K(\xvec, \zvec) = \exp\bigl(- \gamma (\xvec - \zvec)^\t(\xvec - \zvec)\bigr) \text{ with } \gamma > 0. \hfill \inlineeqno$

          If $\gamma = \frac{1}{2\sigma^2}$, then the radial kernel is also called \emph{Gaussian kernel}. For a $\zvec$ far from $\xvec$, the kernel gets small: for this, it is also called \emph{local kernel}.
\end{description}

\begin{example}
    We want to prove that the polynomial kernel of degree $2$ is a valid kernel, \ie, it represents a dot product in some higher-dimensional space. We can write:
    \begin{equation}
        \begin{aligned}
            K(\xvec, \zvec) & = (1 + \xvec^\t\zvec)^2                                                                  \\
                            & = 1 + 2\sum_{j=1}^p x_j z_j + \sum_{j\negthinspace,\,k=1}^p x_jz_j\cdot x_k z_k                         \\
                            & = 1 + \sum_{j=1}^p \sqrt{2}x_j \cdot \sqrt{2}z_j + \sum_{j\negthinspace,\,k=1}^p x_j x_k \cdot z_j z_k.
        \end{aligned}
    \end{equation}
    We can define $\Phi\colon \rr^p \to \rr^{1+p+p^2}$ as:
    \begin{equation}
        \Phi(\xvec) = \tuple{1, \sqrt{2}x_1, \dots, \sqrt{2}x_p, x_1^2, x_1x_2,\ldots, x_p^2} \in \rr^{1+p+p^2}\negmedspace.\qedhere
    \end{equation}
\end{example}

The choice of $\gamma$ and $m$ is important as it relates to the bias-variance trade-off. Again, \acs{cv} is needed for the selection of these parameters.

In conclusion, for an \acs{svm} we need:
\begin{enumerate}
    \item A kernel function $K$ with optimized tuning parameters.
    \item The calculation of $\binom{n}{2}$ values of $K(\xvec_i, \xvec_j)$ to estimate $\beta_0$ and $\gammavec$.
\end{enumerate}

\paragraph{Multiple classes}

If there are more than two classes (namely $K > 2$), we can proceed with mainly two approaches:
\begin{enumerate}
    \item \emph{One vs.\@ All}: train $K$ classifiers where the labels are ``class $c$'' vs.\@ ``not class $c$''\negthinspace, for $c = 1, \dots, K$. One can then combine these classifiers by assigning $\xvec$ to the class $c$ for which $\abs{f_c(\xvec)}$ is maximal.

          However, two problems arise with this approach:
          \begin{itemize}[beginpenalty=10000]
              \item \emph{Potential ambiguity}: if for a given $\xvec$ we have $f_c(\xvec) > 0$ for more than one class, we cannot assign $\xvec$ to a single class.
              \item \emph{Unbalanced data}: if one class is much more represented than the others, the classifier for that class will be more accurate than the others, and it will be more likely to assign $\xvec$ to that class.
          \end{itemize}
    \item \emph{One vs.\@ One}: train $\binom{K}{2} = K(K-1)/2$ classifiers, one for each pair of classes. Then, classify $\xvec$ to the class with the higher number of votes.

          This approach is computationally more expensive, and there could still be potentially ambiguity. However, if $K$ is not too large, this is what is usually done.
\end{enumerate}

\paragraph{\acsp{svc} vs.\@ Logistic Regression}

The primal optimization problem of \acsp{svc} is, as said before:
\begin{equation}
    \begin{aligned}
         & \min_{\beta_0, \betavec, \epsilonvec} \frac{1}{2}\norm{\betavec}^2                                     \\
         & \text{s.t.}\; y_i\left(\beta_0 + \betavec^\t\xvec_i\right) \geq 1 - \epsilon_i, \quad i = 1, \dots, n, \\
         & \phantom{s.t.} \epsilon_i \geq 0, \quad i = 1, \dots, n,                                               \\
         & \phantom{s.t.} \sum_{i=1}^n \epsilon_i \leq C.
    \end{aligned}
\end{equation}
The first constraint is equivalent to:
\begin{equation}
    \epsilon_i \geq 1 - y_i(\beta_0 + \betavec^\t\xvec_i).
\end{equation}
If we combine with the second $\epsilon_i \geq 0$, we get:
\begin{equation}
    \epsilon_i \geq \max\{0, 1 - y_i(\beta_0 + \betavec^\t\xvec_i)\} = [1 - y_i(\beta_0 + \betavec^\t\xvec_i)]_+.
\end{equation}
Using the third constraint:
\begin{equation}
    \sum_{i=1}^n [1 - y_i(\beta_0 + \betavec^\t\xvec_i)]_+ \leq \sum_{i=1}^n \epsilon_i \leq C.
\end{equation}
So, the optimization problem can be rewritten as:
\begin{equation}
    \begin{aligned}
         & \min_{\beta_0, \betavec} \frac{1}{2}\norm{\betavec}^2                        \\
         & \text{s.t.}\; \sum_{i=1}^n [1 - y_i(\beta_0 + \betavec^\t\xvec_i)]_+ \leq C.
    \end{aligned}
\end{equation}
Computing the Lagrangian:
\begin{equation}
    \lagr(\beta_0, \betavec, \gamma) = \frac{1}{2}\norm{\betavec}^2 + \gamma\left(\sum_{i=1}^n [1 - y_i(\beta_0 + \betavec^\t\xvec_i)]_+\right)
\end{equation}
with $\gamma \geq 0$ as a tuning parameter, and the equivalent unconstrained optimization problem is:
\begin{equation}
    \min_{\beta_0, \betavec} \sum_{i=1}^n [1 - y_i(\beta_0 + \betavec^\t\xvec_i)]_+ + \frac{\lambda}{2}\norm{\betavec}^2\negthinspace,
\end{equation}
with $\lambda = 1/\gamma$. Compared with Ridge regression:
\begin{equation}
    \min_{\beta_0, \betavec} \sum_{i=1}^n (y_i - \beta_0 - \betavec^\t\xvec_i)^2 + \frac{\lambda}{2}\norm{\betavec}^2\negthinspace,
\end{equation}
the only difference is the loss function. The loss function of \acsp{svc} is called \emph{hinge loss} and is defined as:
\begin{equation}
    \L(y, f(\xvec)) = \begin{cases}
        0              & \text{if } 1 - y f(\xvec) < 0 \iff yf(\xvec) \geq 1, \\
        1 - y f(\xvec) & \text{if } 1 - y f(\xvec) \geq 0 \iff yf(\xvec) < 1.
    \end{cases}
\end{equation}
We observe that if the point is correctly classified, then the loss is zero. If the point is misclassified, then the loss is proportional to the distance from the decision surface.

\begin{figure}[tb]
    \centering
    \begin{tikzpicture}
        \begin{axis}[
                axis lines = middle,
                xlabel = {$f(\xvec)$},
                ylabel = {$\L(y, f(\xvec))$},
                xtick = {1},
                xticklabels = {$1$},
                ytick = {1},
                yticklabels = {$1$},
                ymax = 1.5,
                xmin = -0.5,
                xmax = 1.75,
                width = 9cm,
                height = 5cm,
            ]
            \addplot[domain=-.35:1.5, samples=5000, thick] {x <= 1 ? 1 - x : 0};
        \end{axis}
    \end{tikzpicture}
    \caption{Hinge loss function.}
\end{figure}

For the logistic regression, we don't have a Ridge penalty, but we can add it. The loss function is the \emph{negative log-likelihood}:
\begin{equation}
    \L(y_i, f(\xvec_i)) = \begin{cases}
        -\log\P(1 \given \xvec_i)       & \text{if } y_i = 1,  \\
        -\log(1 - \P(1 \given \xvec_i)) & \text{if } y_i = -1.
    \end{cases}
\end{equation}
We recall that for logistic regression we have:
\begin{equation}
    \P(1 \given \xvec) = \frac{e^{\beta_0 + \betavec^\t\xvec}}{1 + e^{\beta_0 + \betavec^\t\xvec}}
\end{equation}
and therefore:
\begin{itemize}[beginpenalty=10000]
    \item $y_i = 1$:
          \begin{equation}
              -\log\P(1 \given \xvec_i) = -\log\left(\frac{e^{\beta_0 + \betavec^\t\xvec_i}}{1 + e^{\beta_0 + \betavec^\t\xvec_i}}\right) = \log(1 + e^{-(\beta_0 + \betavec^\t\xvec_i)}).
          \end{equation}
    \item $y_i = -1$:
          \begin{equation}
              -\log(1 - \P(1 \given \xvec_i)) = -\log\left(1 - \frac{e^{\beta_0 + \betavec^\t\xvec_i}}{1 + e^{\beta_0 + \betavec^\t\xvec_i}}\right) = \log(1 + e^{\beta_0 + \betavec^\t\xvec_i}).
          \end{equation}
\end{itemize}
The writing can be simplified by noticing that $y_i \in \{-1, 1\}$:
\begin{equation}
    \L(y_i, f(\xvec_i)) = \log(1 + e^{-y_i(\beta_0 + \betavec^\t\xvec_i)}).
\end{equation}

Let's now compare the two optimization problems:
\begin{description}
    \item[\acs{svc}] $\min_{\beta_0, \betavec} \sum_{i=1}^n [1 - y_i(\beta_0 + \betavec^\t\xvec_i)]_+ + \frac{\lambda}{2}\norm{\betavec}^2\negthinspace$.
    \item[\acs{lr} with Ridge penalty] $\min_{\beta_0, \betavec} \sum_{i=1}^n \log(1 + e^{-y_i(\beta_0 + \betavec^\t\xvec_i)}) + \frac{\lambda}{2}\norm{\betavec}^2\negthinspace$.
\end{description}

Support Vector techniques can be extended to regression problems using the so called \emph{\mbox{$\epsilon$-insensitive} loss function}:
\begin{equation}
    \L(y, f(\xvec)) = \begin{cases}
        0                             & \text{if } \abs{y - f(\xvec)} < \epsilon,    \\
        \abs{y - f(\xvec)} - \epsilon & \text{if } \abs{y - f(\xvec)} \geq \epsilon.
    \end{cases}
\end{equation}
that is zero if the residuals are small enough (less than the parameter $\epsilon$) and proportional to the distance from the decision surface otherwise.

\begin{figure}
    \centering
    \begin{tikzpicture}
        \begin{axis}[
                axis lines = middle,
                xlabel = {$y - f(\xvec)$},
                ylabel = {$\L(y, f(\xvec))$},
                xtick = {-2, -1, 0, 1, 2},
                xticklabels = {$-2$, $-1$, $0$, $1$, $2$},
                ytick = {0, 1},
                yticklabels = {$0$, $1$},
                ymax = 1.5,
                width = 10cm,
                height = 5cm,
            ]
            \addplot[domain=-2.75:2.75, samples=1000, thick] {max(0, abs(x) - 1)};
        \end{axis}
    \end{tikzpicture}
    \caption{$\epsilon$-insensitive loss function with $\epsilon = 1$.}
\end{figure}